% ==============================================================================
% FFGFT: Dokument 247 (Deutsch)
% Titel: Information, Zustand und Prozess — FFGFTs Position
% Autor: Johann Pascher
% Datum: Mai 2026
% ==============================================================================

\section*{Zusammenfassung}

Ist Information Energie? Die Frage wird in der Informationsphysik oft mit Ja
beantwortet. FFGFT differenziert grundlegend:

\begin{itemize}
  \item Die bloße \emph{Speicherung} von Information als Zustand ist energieneutral.
  \item Erst die \emph{Verarbeitung} — insbesondere das \emph{Löschen} — hat
        eine Energiedimension.
\end{itemize}

Diese Unterscheidung ist nicht philosophisch, sondern empirisch entscheidbar:
Permanentspeicher, magnetische Ordnung und der Casimir-Effekt belegen sie
direkt. In FFGFT sind Zustände geometrische Konfigurationen auf $T^4$; Energie
ist eine abgeleitete Größe, die aus Prozessen folgt, nicht aus der bloßen
Existenz von Struktur.

% ------------------------------------------------------------------------------
\section{Die Frage}
% ------------------------------------------------------------------------------

In der Informationsphysik — vertreten etwa durch J.~A.~Wheeler ("It from Bit")
\cite{Wheeler1990} und im IPI-Kontext durch Denis \cite{Denis2023} — wird
Information oft als ontologisches Primitiv behandelt, aus dem Energie und
Materie folgen. Landauers Prinzip \cite{Landauer1961} wird dabei häufig als
Beleg für die physikalische Realität von Information angeführt: da das Löschen
eines Bits Energie kostet, sei Information physikalisch.

Dieser Schluss verwechselt zwei strukturell verschiedene Begriffe:

\begin{keyresult}[Zentrale Unterscheidung]
\textbf{Information als Zustand} — eine Konfiguration, ein gespeichertes Bit,
eine Wicklungszahl — existiert ohne Energieaufwand. Sie ist eine topologische
oder geometrische Eigenschaft.\\
\textbf{Information als Prozess} — Lesen, Schreiben, Löschen, Transformieren
— kostet Energie. Landauers Prinzip gilt \emph{nur} für die Löschung, also
für einen \emph{Prozess}, nicht für den Zustand selbst.
\end{keyresult}

Wer diesen Unterschied einebnet, muss eine direkte empirische Konsequenz
akzeptieren: Jedes gespeicherte Bit auf einem ausgeschalteten Gerät würde
kontinuierlich Energie verbrauchen — was Permanentspeicher grundsätzlich
unmöglich machen würde. Die Existenz solcher Speicher widerlegt die
Gleichsetzung von Information und Energie.

% ------------------------------------------------------------------------------
\section{Landauers Prinzip präzise}
% ------------------------------------------------------------------------------

Rolf Landauer zeigte 1961, dass das irreversible Löschen eines Bits
Wärme an die Umgebung abgibt \cite{Landauer1961}. Die Mindestenergie beträgt:

\begin{equation}
  E_{\min} = k_B T \ln 2
\end{equation}

wobei $k_B$ die Boltzmann-Konstante und $T$ die absolute Temperatur ist.
Bei Raumtemperatur entspricht das etwa $2{,}9 \times 10^{-21}\,\text{J}$
pro gelöschtem Bit.

\textbf{Was Landauer zeigt:} Die Löschung — ein irreversibler Prozess —
dissipiert Energie. Das ist eine Konsequenz des zweiten Hauptsatzes der
Thermodynamik \cite{Bennett1982}.

\textbf{Was Landauer nicht zeigt:} Dass die bloße \emph{Existenz} eines
gespeicherten Bits Energie erfordert. Ein magnetisches Bit auf einer
Festplatte, eine Windungszahl auf $T^4$, eine Spinausrichtung im Permanentmagneten
— all das existiert ohne laufenden Energiebedarf. Erst der Übergang zwischen
Zuständen kostet.

In natürlichen Einheiten ($\hbar = c = k_B = 1$) hat die Löschungsenergie
die Dimension einer Temperatur, also einer Masse, also einer Energie — und
damit Anschluss an $\tilde{T} \cdot m = 1$ in FFGFT. Aber diese
Dimensionsbeziehung gilt für den \emph{Prozess}, nicht für den ruhenden
Zustand.

% ------------------------------------------------------------------------------
\section{Einschreiben versus Löschen: kein symmetrischer Energiebedarf}

Man könnte einwenden: auch das Einschreiben eines Bits kostet Energie —
und wenn Energie eingebracht wird, entsteht nach $E = mc^2$ Masse. Wäre
Information also doch massebehaftet?

Dieser Schluss ist aus zwei Gründen falsch.

\subsection{Logische Reversibilität entscheidet, nicht der Prozess als solcher}

Die vier möglichen deterministischen Operationen auf einem Bit sind:
HOLD (Halten), RESET-to-0, RESET-to-1 und NOT (Flip). Davon sind
RESET-to-0 und RESET-to-1 logisch \emph{irreversibel} — der vorherige
Zustand geht verloren. NOT und HOLD sind logisch \emph{reversibel} —
der vorherige Zustand ist aus dem Ergebnis rekonstruierbar.

Landauers Schranke gilt \emph{nur} für logisch irreversible Operationen,
also nur für RESET \cite{Landauer1961}. Dago und Bellon zeigten
experimentell, dass die NOT-Operation — ein Einschreiben eines neuen
Zustands durch Bitflip — in einem physikalisch reversiblen Verfahren
mit verschwindend geringen Energiekosten durchgeführt werden kann, weit
unterhalb der Landauer-Schranke \cite{Dago2023}. Es gibt keine
fundamentale Energieschranke für das Einschreiben an sich.

\subsection{Energie fließt in die Umgebung, nicht in den Zustand}

Praktisch kostet das Einschreiben immer Energie — aber diese Energie
geht nicht \emph{in das gespeicherte Bit}. Sie geht in:

\begin{itemize}
  \item Ohmsche Wärme im Schreibstrom (Flash, DRAM)
  \item Magnetische Ummagnetisierungsarbeit gegen die Koerzitivkraft
  \item Mechanische Dissipation im Schreibkopf
\end{itemize}

Das gespeicherte Bit selbst trägt keine zusätzliche Energie oder Masse.
Ein beschriebener Flash-Speicher wiegt nicht mehr als ein leerer —
obwohl das Schreiben Energie gekostet hat. Die Energie ist in die
Umgebung abgeflossen, nicht im Zustand gespeichert.

\subsection{Warum entsteht keine Masse durch Information}

Der Schluss „Einschreiben kostet Energie → $E=mc^2$ → Masse entsteht"
begeht einen Kategorienfehler. $E=mc^2$ gilt für die Ruheenergie eines
physikalischen Systems — also für die Masse des Speichers selbst, nicht
für die Energie die beim Schreibprozess in die Umgebung dissipiert wird.

In FFGFT ist das konsistent: Wicklungszahlen auf $T^4$ sind topologische
Invarianten. Der Übergang von einer Konfiguration zur anderen ist ein
Prozess mit Energiedimension über $\tilde{T}\cdot m=1$ — aber die
Konfiguration selbst trägt keine zusätzliche Ruhemasse. Der Prozess
erzeugt keine neue Masse, er verändert den Zustand eines Systems das
bereits Masse hat.

\section{Drei empirische Pfeiler für die Energieneutralität von Zuständen}

\subsection{Pfeiler 1: Magnetismus}

Ein Permanentmagnet ist der anschaulichste Beleg für die Energieneutralität
von Informationszuständen. Die Spins in einem ferromagnetischen Material sind
ausgerichtet — das ist ein Informationszustand im physikalischen Sinne
(Nordpol vs.\ Südpol, Bit = Magnetisierungsrichtung).

\begin{itemize}
  \item \textbf{Speicherung:} kostet keine laufende Energie. Der Magnet
        hält seine Ausrichtung ohne Stromzufuhr — prinzipiell unbegrenzt.
  \item \textbf{Ummagnetisierung:} kostet Energie. Das Überwinden der
        Koerzitivfeldstärke erfordert ein externes Feld, das Arbeit verrichtet.
  \item \textbf{Langsamer Zerfall:} Über sehr lange Zeiträume können
        thermische Fluktuationen, kosmische Strahlung oder andere externe
        Störungen die Ausrichtung ändern \cite{Nattermann1990}. Aber dieser
        Zerfall ist extern ausgelöst — ein von außen getriggerter Prozess,
        keine intrinsische Eigenschaft des Zustands selbst.
\end{itemize}

Das ist der strukturelle Punkt: selbst der langsame Zerfall eines
Permanentspeichers ist ein Prozess mit externer Ursache, kein
Beweis dafür, dass der Zustand selbst Energie kostet.

\subsection{Pfeiler 2: Casimir-Effekt}

Der Casimir-Effekt zeigt dass Vakuumgeometrie physikalische Konsequenzen
hat, ohne dass ein Energiespeicher im klassischen Sinne vorhanden ist.
Zwei ungeladene, parallele Metallplatten im Vakuum ziehen einander an —
eine messbare Kraft, die 1948 von Hendrik Casimir vorhergesagt wurde
\cite{Casimir1948} und 1996 von Steve Lamoreaux erstmals präzise gemessen
wurde \cite{Lamoreaux1997}.

Die Kraft pro Fläche beträgt:

\begin{equation}
  P(z) = -\frac{\pi^2}{240}\frac{\hbar c}{z^4}
\end{equation}

wobei $z$ der Plattenabstand ist.

\textbf{Was der Casimir-Effekt zeigt:} Die geometrische Einschränkung der
erlaubten Vakuummoden zwischen den Platten erzeugt einen Energiedichteunterschied,
der eine Kraft bewirkt. Die \emph{Geometrie des Vakuums} — der Zustand — hat
physikalische Konsequenzen.

\textbf{Was der Casimir-Effekt nicht zeigt:} Dass diese Geometrie Energie
speichert oder verbraucht, solange die Platten stillstehen. Erst die
\emph{Bewegung der Platten} — der Prozess — verrichtet oder nimmt Arbeit auf.

Die strukturelle Parallele zu FFGFT ist folgende: Die Wicklungsstruktur
auf $T^4$ ist eine geometrische Konfiguration — sie existiert, hat
Konsequenzen (Massenverhältnisse, Kopplungskonstanten), aber sie
\emph{kostet} keine laufende Energie. Erst dynamische Übergänge zwischen
Konfigurationen — Prozesse — haben Energiedimension. Diese Parallele ist
eine Strukturanalogie, keine Identität der Mechanismen.

\subsection{Pfeiler 3: Permanentspeicher}

Die gesamte digitale Zivilisation basiert auf der Energieneutralität
gespeicherter Information. Flash-Speicher, magnetische Festplatten,
optische Datenträger — alle nutzen Zustände die ohne Stromzufuhr bestehen.

\begin{itemize}
  \item \textbf{Flash-Speicher} speichert Bits als Ladungszustände in
        Floating-Gate-Transistoren. Die Ladung bleibt ohne Energiezufuhr
        erhalten — typischerweise 10 Jahre oder länger.
  \item \textbf{Magnetbänder} halten Daten über Jahrzehnte ohne
        Stromversorgung.
  \item \textbf{Langsamer Zerfall:} Quantentunneling, thermische
        Aktivierung und ionisierende Strahlung können Bits über sehr
        lange Zeiträume kippen \cite{Bez2003}. Aber auch hier: jedes
        einzelne Kippen ist ein von außen getriggerter Prozess, keine
        spontane Eigenschaft des Ruhezustands.
\end{itemize}

Wäre Information als Zustand nicht energieneutral, wäre kein Permanentspeicher
möglich. Die Existenz dieser Technologie ist der empirische Beweis für die
zentrale Unterscheidung.

% ------------------------------------------------------------------------------
\section{In natürlichen Einheiten: Zustand, Prozess und die Grundrelation}
% ------------------------------------------------------------------------------

In natürlichen Einheiten ($\hbar = c = k_B = 1$) kollabieren die
Dimensionen:

\begin{equation}
  [E] = [m] = [T^{-1}] = [\text{Temperatur}]
\end{equation}

Die Löschungsenergie $E_{\min} = k_B T \ln 2$ wird in natürlichen Einheiten
zu $E_{\min} = T \ln 2$ — eine Temperatur, eine Masse, eine Energie.

FFGFTs Grundrelation $\tilde{T} \cdot m = 1$ verbindet Zeit und Masse direkt.
Ein Zustandsübergang — ein Prozess — hat Zeitdauer und verändert die effektive
Masse des dynamischen Feldes $\Delta m(x,t)$. Das ist die natürliche Einbettung
von Landauers Prinzip in FFGFT:

\begin{itemize}
  \item \textbf{Zustand:} Wicklungskonfiguration auf $T^4$ — $m$ ist
        geometrisch fixiert durch $\xi$, kein Energieverbrauch.
  \item \textbf{Prozess:} Übergang zwischen Konfigurationen — $\Delta m$
        ändert sich, das hat über $\tilde{T} \cdot m = 1$ eine
        Zeitdimension und damit eine Energiedimension.
\end{itemize}

Die Energiekosten von Informationsverarbeitung sind in FFGFT also nicht
als externen Zusatz einzuführen — sie folgen aus der Grundrelation.

% ------------------------------------------------------------------------------
\section{FFGFTs Position}
% ------------------------------------------------------------------------------

\begin{keyresult}[FFGFT und Information]
\textbf{Primitiv:} Geometrie — $T^4$-Wicklungsstruktur, $\xi$, $\tilde{T}\cdot m=1$.\\
\textbf{Zustand:} Wicklungskonfiguration — kein Energieprimitivum, analog
zum magnetischen Bit oder zur Casimir-Geometrie.\\
\textbf{Prozess:} Zustandsübergang — hat Energiedimension, folgt aus
$\tilde{T}\cdot m=1$ und $\xi$.\\
\textbf{Information:} Beschreibungsebene über der geometrischen Realität,
kein ontologisches Primitiv.
\end{keyresult}

Shannon-Information \cite{Shannon1948} — Bits, Entropie $H = -\sum p_j \log p_j$
— ist eine Zählgröße, dimensionslos. Sie beschreibt Zustände in einem
Wahrscheinlichkeitsraum. Sie ist kein Primitiv in FFGFT, weil FFGFT
keine Wahrscheinlichkeiten als Grundbegriff verwendet — die Wicklungszahlen
sind deterministische topologische Invarianten, keine Zufallsvariablen.

% ------------------------------------------------------------------------------
\section{Abgrenzung}
% ------------------------------------------------------------------------------

\subsection{Wheeler — „It from Bit"}

John Archibald Wheelers Position \cite{Wheeler1990}: Information ist das
ontologische Primitiv, aus dem physikalische Realität entsteht. Das setzt
Information vor Geometrie.

FFGFT kehrt das um: Geometrie ($T^4$, $\xi$) ist das Primitiv. Information
im Shannon-Sinne ist eine Beschreibung von Konfigurationen in dieser Geometrie,
keine Ursache davon.

\subsection{Denis — Landauer-Vakuum (IPI)}

Denis \cite{Denis2023} schlägt vor, die kosmologische Konstante über
Landauers Prinzip aus einem informationstheoretischen Vakuumsektor
abzuleiten. Das setzt voraus, dass Information ontologisch fundamental ist.

FFGFT leitet $\Lambda = 1{,}34 \times 10^{-122}$ (Planck-Einheiten)
direkt aus $\xi$ ab \cite{Pascher_182}, ohne Informationstheorie als
Zwischenschicht. Der Weg ist kürzer und benötigt keine Landauer-Brücke.

\subsection{RSG — Survival-Gewichte}

RSG \cite{Austin2026} verwendet Überlebensgewichte $S_i = e^{-A_i}$ für
Historien. Das ist eine Wahrscheinlichkeitsstruktur über Konfigurationen —
also eine Beschreibungsebene, kein geometrisches Primitiv. Aus FFGFT-Sicht
ist RSG eine mögliche Beschreibung von Selektionsprozessen, aber keine
Erklärung der Skalenwerte die $\Gamma_i$ bestimmen.

% ------------------------------------------------------------------------------
\section{Zeichenerklärung}
% ------------------------------------------------------------------------------

{\small
\adjustbox{max width=\textwidth}{%
\begin{tabular}{p{3.5cm}p{9.5cm}}
\toprule
\textbf{Symbol} & \textbf{Bedeutung} \\
\midrule
$k_B$ & Boltzmann-Konstante: $k_B = 1{,}380649 \times 10^{-23}\,\text{J/K}$ \\
$T$ & Absolute Temperatur (Kelvin) \\
$E_{\min} = k_B T \ln 2$ & Landauer-Schranke: Mindestenergie zum Löschen eines Bits \\
$H = -\sum p_j \log p_j$ & Shannon-Entropie; dimensionslos in Bits (Basis 2) \cite{Shannon1948} \\
$\xi = 4/30000$ & FFGFT-Kopplungsparameter; dimensionslos \\
$\tilde{T} \cdot m = 1$ & FFGFT-Grundrelation (natürliche Einheiten) \\
$T^4$ & 4-dimensionaler Torus; ontologisches Primitiv von FFGFT \\
$\Delta m(x,t)$ & Dynamisches Massenfeld in T0; trägt Zustandsübergänge \\
$P(z) = -\frac{\pi^2}{240}\frac{\hbar c}{z^4}$ & Casimir-Druck; $z$ = Plattenabstand \cite{Casimir1948} \\
$\Lambda = 1{,}34 \times 10^{-122}$ & Kosmologische Konstante aus $\xi$ (Planck-Einheiten) \cite{Pascher_182} \\
\bottomrule
\end{tabular}%
}
}

% ------------------------------------------------------------------------------
\section*{Literatur}
% ------------------------------------------------------------------------------

\begin{thebibliography}{99}

\bibitem{Landauer1961}
R.~Landauer,
\textit{Irreversibility and Heat Generation in the Computing Process},
IBM Journal of Research and Development \textbf{5}(3), 183--191 (1961).
DOI: 10.1147/rd.53.0183.

\bibitem{Bennett1982}
C.~H.~Bennett,
\textit{The Thermodynamics of Computation — A Review},
International Journal of Theoretical Physics \textbf{21}(12), 905--940 (1982).
DOI: 10.1007/BF02084158.

\bibitem{Shannon1948}
C.~E.~Shannon,
\textit{A Mathematical Theory of Communication},
Bell System Technical Journal \textbf{27}(3), 379--423 (1948).
DOI: 10.1002/j.1538-7305.1948.tb01338.x.

\bibitem{Casimir1948}
H.~B.~G.~Casimir,
\textit{On the Attraction Between Two Perfectly Conducting Plates},
Proceedings of the Koninklijke Nederlandse Akademie van Wetenschappen
\textbf{51}, 793--795 (1948).

\bibitem{Lamoreaux1997}
S.~K.~Lamoreaux,
\textit{Demonstration of the Casimir Force in the 0.6 to 6 $\mu$m Range},
Physical Review Letters \textbf{78}(1), 5--8 (1997).
DOI: 10.1103/PhysRevLett.78.5.

\bibitem{Wheeler1990}
J.~A.~Wheeler,
\textit{Information, Physics, Quantum: The Search for Links},
in: W.~H.~Zurek (Hrsg.), \textit{Complexity, Entropy, and the Physics of
Information}, Addison-Wesley, Redwood City (1990), S.~3--28.

\bibitem{Nattermann1990}
T.~Nattermann,
\textit{Theory of the Random Field Ising Model},
in: A.~P.~Young (Hrsg.), \textit{Spin Glasses and Random Fields},
World Scientific, Singapur (1998), S.~277--298.

\bibitem{Bez2003}
R.~Bez et al.,
\textit{Introduction to Flash Memory},
Proceedings of the IEEE \textbf{91}(4), 489--502 (2003).
DOI: 10.1109/JPROC.2003.811702.

\bibitem{Denis2023}
O.~Denis,
\textit{Informational Nature of Dark Matter and Dark Energy and the
Cosmological Constant},
IPI Letters \textbf{1}, 66--77 (2023).
DOI: 10.59973/ipil.36.

\bibitem{Austin2026}
P.~M.~Austin,
\textit{Recursive Survival Geometry: A Conditional Recursive Formalism
and Analogue-Media Test Protocol for Survival-Weighted Representation},
v1.06, Information Physics Institute (2026).

\bibitem{Dago2023}
S.~Dago und L.~Bellon,
\textit{Logical and Thermodynamical Reversibility: Optimized Experimental
Implementation of the NOT Operation},
arXiv:2302.11908 (2023).

\bibitem{Pascher_182}
J.~Pascher,
\textit{T0-Universum Maximalskala — Kosmologische Konstante und Hubble-Parameter
aus $\xi$},
Dokument 182, FFGFT-Dokumentenserie (2026).
GitHub: jpascher/T0-Time-Mass-Duality.
Zenodo: DOI 10.5281/zenodo.20474821.

\end{thebibliography}
